\chapter{Statement of the Problem}
\label{chap:problem_statement}

\section{Problem Statement}
\label{sec:problem_statement}

Traffic congestion poses a significant barrier to a country’s economic growth, leading to environmental pollution and a decline in the quality of life (Fattah, 2022). In the Philippines, rapidly growing urban centers, including Iloilo City, have driven increased demand for transportation. According to (Landoy, 2012), as cited in (Sosuan, 2014), traffic congestion is already eminent along Iloilo City’s primary roads. A common intuitive solution to congestion is to expand capacity by building new roads or adding links. However, empirical evidence from decades of traffic data across the United States shows that increasing road capacity does not necessarily alleviate congestion (Stromberg, 2014). 

This global trend demonstrates that similar outcomes may also occur in developing urban centers such as Iloilo City. This counterintuitive phenomenon, known as the Braess Paradox, describes a situation in which strategically adding a new road link results in an overall increase in travel costs for all network users, rather than reducing them (Rapoport, 2009). The potential for such an outcome in Iloilo City highlights a critical gap between intuitive infrastructure planning and actual traffic network behavior. Thus, there is a pressing need for an agent-based modeling (ABM) framework that can identify such paradoxical outcomes in a local setting. This study addresses this need by developing and testing a model within a focused study area, which serves as the foundational case study for this analysis.